% Part: first-order-logic
% Chapter: tableaux
% Section: quantifier-rules

\documentclass[../../../include/open-logic-section]{subfiles}

\begin{document}

\olfileid{fol}{tab}{qrl}

\olsection{量词规则}

\subsection{$\lforall$ 的规则}

\begin{defish}
\AxiomC{\sFmla{\True}{\lforall[x][!A(x)]}}
\RightLabel{\TRule{\True}{\forall}}
\UnaryInfC{\sFmla{\True}{!A(t)}}
\DisplayProof
\hfill
\AxiomC{\sFmla{\False}{\lforall[x][!A(x)]}}
\RightLabel{\TRule{\False}{\lforall}}
\UnaryInfC{\sFmla{\False}{!A(a)}}
\DisplayProof
\end{defish}

在 \TRule{\True}{\lforall} 规则中，$t$ 是一个闭项（即不含变元的项）。在 \TRule{\False}{\lforall} 规则中，$a$~是一个常项符号，且它不得出现在 \TRule{\False}{\lforall} 规则上方的分支中的任何地方。我们称 $a$ 为 \TRule{\False}{\lforall} 推理的\emph{本征变元}。\footnote{尽管上述规则中的 $a$ 是常项符号，我们仍沿用“本征变元”这个术语。这是出于历史原因。}

\subsection{$\lexists$ 的规则}

\begin{defish}
\AxiomC{\sFmla{\True}{\lexists[x][!A(x)]}}
\RightLabel{\TRule{\True}{\lexists}}
\UnaryInfC{\sFmla{\True}{!A(a)}}
\DisplayProof
\hfill
\AxiomC{\sFmla{\False}{\lexists[x][!A(x)]}}
\RightLabel{\TRule{\False}{\lexists}}
\UnaryInfC{\sFmla{\False}{!A(t)}}
\DisplayProof
\end{defish}

同样地，$t$~是闭项，而 $a$~是一个常项符号，且不得出现在 \TRule{\True}{\lexists} 规则上方的分支中。我们称 $a$ 为 \TRule{\True}{\lexists} 推理的\emph{本征变元}。

本征变元不得在 \TRule{\False}{\lforall} 或 \TRule{\True}{\lexists} 推理上方的分支中出现，这一条件被称为\emph{本征变元条件}。

\begin{explain}
回顾一下记法约定：当 $!A$ 是一个变元~$x$ 在其中自由出现的公式时，我们写作~$!A(x)$。在同一语境下，$!A(t)$ 则是~$\Subst{!A}{t}{x}$ 的简写。因此，我们也可以将 \TRule{\False}{\lexists} 规则写作：
\begin{prooftree}
  \AxiomC{\sFmla{\False}{\lexists[x][!A]}}
  \RightLabel{\TRule{\False}{\lexists}}
  \UnaryInfC{\sFmla{\False}{\Subst{!A}{t}{x}}}
\end{prooftree}
注意，$t$ 可能已经在~$!A$ 中出现，例如，$!A$~可能是~$\Atom{\Obj P}{t,x}$。因此，从~$ \sFmla{\False}{\lexists[x][\Atom{\Obj P}{t,x}]}$ 推出 $\sFmla{\False}{\Atom{\Obj
P}{t,t}}$ 是~\TRule{\False}{\lexists} 的正确应用。但是，\TRule{\False}{\lforall} 和~\TRule{\True}{\lexists} 中的本征变元条件要求常项符号~$a$ 不得出现在~$!A$ 中。所以，你无法使用~$\TRule{\False}{\lforall}$ 从
$\sFmla{\False}{\lforall[x][\Atom{\Obj P}{a,x}]}$
正确地推出
$\sFmla{\False}{\Atom{\Obj P}{a,a}}$。
\end{explain}

\begin{explain}
在 \TRule{\True}{\lforall} 和 \TRule{\False}{\lexists} 中，对项~$t$ 没有限制。另一方面，在 \TRule{\True}{\lexists} 和 \TRule{\False}{\lforall} 规则中，本征变元条件要求常项符号~$a$ 不得出现在各自推理上方分支的任何地方。这一条件对于保证系统的可靠性是必要的。如果没有这个条件，下面就会是 $\lexists[x][\formula{A}(x)] \lif \lforall[x][\formula{A}(x)]$ 的一个封闭分析表：
\begin{center}
\begin{tableau}{}
  [\sFmla{\False}{\lexists[x][\formula{A}(x)] \lif \lforall[x][\formula{A}(x)]}, just=\TAss
    [\sFmla{\True}{\lexists[x][\formula{A}(x)]},
      just={\TRule{\False}{\lif}[1]}
      [\sFmla{\False}{\lforall[x][\formula{A}(x)]},
        just={\TRule{\False}{\lif}[1]}
        [\sFmla{\True}{\formula{A}(a)},
          just={\TRule{\True}{\lexists}[2]}
          [\sFmla{\False}{\formula{A}(a)},
            just={\TRule{\False}{\lforall}[3]}, close]
        ]
      ]
    ]
  ]
\end{tableau}
\end{center}
然而，$\lexists[x][\formula{A}(x)] \lif
\lforall[x][\formula{A}(x)]$ 并非有效式。
\end{explain}

\end{document}